\documentclass{reportkit}
\setreportkitleftheader{REPORTKIT / ALGORITHM P3 TEST}
\setreportkitfooter{P3 algorithm-state extensions}
\title{ReportKit -- P3 algorithm-state extensions}
\author{ReportKit acceptance fixture}

\begin{document}
\maketitle

\section{Keyed join}
\begin{diagram}[type=join,caption={Hash join state.},source={Conceptual algorithm visual.},description={Orders and customers feed a central hash index. Customer C42 matches and produces a joined row while customer C99 is discarded as unmatched.}]
  \begin{joinstate}
    \joininput{left}{Orders}
    \joininput{right}{Customers}
    \hashbucket[state=active]{customer\_id}{C42,C73}
    \joinmatch{customer\_id}{C42 \(\to\) order/customer row}
    \joinunmatched{right}{C99}
    \joinoutput{Joined rows}
  \end{joinstate}
\end{diagram}

\section{Disjoint sets}
\begin{diagram}[type=union-find,caption={Union/find state.},source={Conceptual algorithm visual.},description={Four nodes are laid out as a disjoint-set forest. B points to representative A, a union operation joins A and C, and a find path runs from D toward A.}]
  \begin{unionfindstate}[columns=4]
    \ufnode{A}{A}
    \ufnode{B}{B}
    \ufnode{C}{C}
    \ufnode[state=current]{D}{D}
    \parent{B}{A}
    \union{A}{C}
    \findpath{D}{A}
  \end{unionfindstate}
\end{diagram}

\section{Linked list}
\begin{diagram}[type=linked-list,caption={Next-pointer traversal.},source={Conceptual algorithm visual.},description={A linked list runs from HEAD through a current middle node to TAIL and then to an explicit NULL terminator.}]
  \begin{linkedliststate}
    \listnode{a}{A}
    \listnode[state=current]{b}{B}
    \listnode{c}{C}
    \nextlink{a}{b}
    \nextlink{b}{c}
    \head{a}
    \tail{c}
  \end{linkedliststate}
\end{diagram}

\section{Recursion}
\begin{diagram}[type=recursion,caption={Memoized recursion.},source={Conceptual algorithm visual.},description={A recursive call view shows Fibonacci arguments and return values. One subproblem is marked MEMOIZED and the root call is current.}]
  \begin{recursiontree}[columns=3]
    \recursionnode[arguments={n=4},state=current]{f4}{fib}{3}
    \recursionnode[arguments={n=3}]{f3}{fib}{2}
    \recursionnode[arguments={n=2},memoized]{f2}{fib}{1}
    \recursionedge{f4}{f3}
    \recursionedge{f3}{f2}
  \end{recursiontree}
\end{diagram}

\end{document}
